\section{Related Work}\label{sec:related}

The development of quantum programming languages and frameworks has evolved significantly over the past two decades, with researchers proposing various approaches to address the challenges of quantum-classical integration, abstraction levels, and error handling. This section systematically reviews significant contributions to quantum programming models and frameworks, focusing particularly on their approaches to managing the hybrid nature of practical quantum computation, which is central to our work.

\subsection{Quantum Programming Languages}

The evolution of quantum programming languages reflects the ongoing effort to balance theoretical rigor with practical implementation considerations. We group these languages based on their design philosophies and implementation strategies, recognizing that these categories are not strictly exclusive, as many languages combine elements from multiple approaches.

\textbf{Foundational and Theoretical Languages} have established key theoretical principles for quantum programming while often incorporating high-level abstractions. Peter Selinger's QPL~\cite{selinger2004towards} (2004) introduced rigorous denotational semantics using superoperators on density matrices but lacked practical hardware implementation. Quipper~\cite{green2013quipper} (2013) leveraged Haskell's functional paradigm as a standalone DSL, introducing high-level quantum data types and circuit transformations as first-class objects while relying on host language mechanisms for classical control. Twist~\cite{yuan2022twist} (2022) addressed entanglement management through a sophisticated type system for quantum state purity. Despite their different focuses and high-level features, all maintain a model where classical computation directs quantum operations—a limitation our stack-based callback model specifically addresses.

\textbf{Quantum Assembly Languages} serve as intermediaries between high-level programming and hardware execution. eQASM~\cite{fu2019eqasm} separated quantum operations from timing specifications with a focus on superconducting processors. cQASM~\cite{khammassi2018cqasm} aimed to standardize quantum circuit representation across platforms. QUASAR~\cite{butko2020understanding}, extending RISC-V, focused on efficient circuit encoding and real-time requirements. These approaches offer low-level hybrid programming models supporting both classical and quantum instructions on a single processor. However, programming with sophisticated classical feedback requires extensive boilerplate code and lacks the high-level composability that qstack provides through its callback model.

\textbf{Hardware-Oriented Languages} prioritize abstractions that match existing quantum hardware capabilities. OpenQASM~\cite{cross2017open,cross2022openqasm}, widely adopted in the IBM ecosystem, includes classical control flow, user-defined gates, and timing constraints, though with rudimentary classical processing capabilities. Quil~\cite{smith2016practical} features ``classical memory regions'' for quantum-classical communication but with limited expressiveness. These approaches differ from our stack-based model in two ways: they rely on explicitly named registers rather than dynamically managed kernel structures, and their classical capabilities lack the expressiveness needed for sophisticated error correction systems. This limits their application to practical error correction schemes requiring complex decoding algorithms or adaptive measurement sequences.

\textbf{Integrated Quantum-Classical Languages} focus on making quantum algorithm development more accessible with native support for both paradigms. Microsoft's Q\#~\cite{qsharp} and Silq~\cite{silq} include built-in classical programming constructs rather than relying on host languages. Q\# implements strong type safety and specialized quantum data types, while Silq innovates with automatic uncomputation and sophisticated type systems. Qunity~\cite{voichick2023qunity} demonstrates the significant challenges inherent in formally combining classical and quantum programming models, requiring complex semantics and type systems that must bridge fundamentally different computational paradigms. These and other languages (including aspects of Quipper and Twist) represent ongoing efforts to create accessible quantum programming environments. While they provide elegant abstractions for algorithm development, many struggle to integrate advanced classical algorithms needed for sophisticated error correction decoding. Our stack-based model treats quantum execution as primary while enabling seamless integration with full-featured classical programming through callbacks.

\subsection{Quantum Error Correction Frameworks}

Specialized frameworks for quantum error correction bridge theoretical protocols and practical implementation. Stim~\cite{gidney2021stim}, developed at Google Quantum AI, represents the state-of-the-art in error correction simulation, efficiently handling stabilizer circuits with tens of thousands of qubits through optimized tableau representations. It provides classical feedback channels, detector-based error analysis, and efficient sampling for characterizing error correction performance under realistic noise. Stim's approach to classical feedback aligns conceptually with our work, though implemented through a specialized engine for stabilizer circuits rather than a general-purpose model. While Stim excels at simulating error correction protocols, qstack integrates error correction within a broader quantum programming framework.


\subsection{Classical-Quantum Integration Models}

Classical-quantum integration presents significant challenges, with several prevailing approaches: (1) \textbf{Host language embedding} (Qiskit~\cite{Qiskit}, Cirq~\cite{cirq}, PyQuil~\cite{pyquil}) embeds quantum operations within classical languages, offering flexibility but obscuring runtime requirements for complex quantum programs. XACC~\cite{mccaskey2020xacc} extends this approach as a C++ embedded DSL with a hardware-agnostic service architecture for heterogeneous programming model integration; (2) \textbf{Extended circuit models} (OpenQASM 3.0, Quil) add classical control to circuit descriptions but limit expressiveness for complex measurement processing; (3) \textbf{Specialized hybrid languages} like Q\#~\cite{qsharp} unify quantum-classical constructs but lack expressiveness for complex decoding algorithms like Minimum Weight Perfect Matching~\cite{higgott2022pymatching}. Despite these varied approaches, none treat classical components as opaque callbacks within a primarily quantum program, even foundational approaches maintain classical computation as the driver of quantum operations. Our stack-based approach inverts this relationship—treating quantum programs as primary with classical operations integrated as callbacks—simplifying algorithms requiring intricate measurement and feedback patterns. This opaqueness of classical components is what enables qstack to be entirely agnostic about which classical language is used or the complexity of classical processing required, from simple bit manipulation to sophisticated error-decoding algorithms.